\chapter{力与运动}\label{ch:g11:forces-and-motion}

将一把旋转着的扳手抛过房间：两端剧烈晃动，然而
其内部某一点却画出一条完美的抛物线。\Cref{ch:g10:inertia}
已经说明：各力相互抵消时——什么也不会发生。
本章则探讨各力不能相互抵消时会发生什么：剩余的力\emph{改变
速度}——使物体加速、减速或转弯，
——并说明改变了多少，此时只用比例，尚未写成公式。

\section{质心}

\begin{definition}[质心]\label{def:g11:forces-and-motion:com}
物体或系统的\emph{质心}\index{质心}是
它的平衡点——即质量的平均位置：
对于质量均匀的对称物体为其几何中心；对于两个质点，质心位于
两者连线上，到两质点的距离与质量成\emph{反}比，
$\frac{d_1}{d_2} = \frac{m_2}{m_1}$——更靠近质量较大的一端。
\end{definition}

\begin{example}[不对称的哑铃]\label{ex:g11:forces-and-motion:dumbbell}
质量分别为 \qty{2.0}{kg} 和 \qty{1.0}{kg} 的小球，固定在一根
长 \qty{0.90}{m} 的轻杆两端。质心按反比分割杆长，
$d_1/d_2 = 1.0/2.0$：距重球 $\qty{0.30}{m}$，
距轻球 $\qty{0.60}{m}$——在此处支撑杆即可使其平衡。
\end{example}

\begin{remark}[抛出的扳手]\label{rem:g11:forces-and-motion:wrench}
拍摄翻滚中的扳手：每一点都沿着复杂的环形曲线运动
——唯有质心沿着被抛出小球会走的那条平滑抛物线运动。
质心的运动很简单；其余部分
是\emph{绕}质心的自转——这正是力学中可以把汽车、行星或
体操运动员视为一个质点的原因。
\end{remark}

\begin{omfigure}
\begin{tikzpicture}[scale=1.0, font=\small]
  \draw[gray!60] (-0.4,0) -- (8.6,0);
  \draw[dashed, omDef, thick]
    plot[domain=0:8.2, samples=60] (\x, {0.9*\x - 0.11*\x*\x});
  \draw[very thick, omExo] ($(0.5,0.42)+(20:0.55)$) -- ($(0.5,0.42)-(20:0.55)$);
  \draw[very thick, omExo] ($(2.3,1.49)+(95:0.55)$) -- ($(2.3,1.49)-(95:0.55)$);
  \draw[very thick, omExo] ($(4.1,1.84)+(170:0.55)$) -- ($(4.1,1.84)-(170:0.55)$);
  \draw[very thick, omExo] ($(5.9,1.48)+(245:0.55)$) -- ($(5.9,1.48)-(245:0.55)$);
  \draw[very thick, omExo] ($(7.7,0.41)+(320:0.55)$) -- ($(7.7,0.41)-(320:0.55)$);
  \fill[omDef] (0.5,0.42) circle (2pt) (2.3,1.49) circle (2pt)
    (4.1,1.84) circle (2pt) (5.9,1.48) circle (2pt) (7.7,0.41) circle (2pt);
  \node[omDef] at (4.1,2.4) {质心};
\end{tikzpicture}

{\small 相等时间间隔下被抛出的扳手：扳手在自转，但
其质心（圆点）沿一条平滑抛物线运动。}
\end{omfigure}

\section{合力改变速度}

\begin{definition}[合力]\label{def:g11:forces-and-motion:netforce}
物体所受的\emph{合力}\index{合力}是
作用在其上的所有力的矢量和，$\vect F = \vect F_1 + \vect F_2 + \dots$
（首尾相接，如数学分册所述）；
当 $\vect F = \vect 0$ 时，各力恰好相互抵消。
\end{definition}

\begin{theorem}[合力的作用效果——半定量]\label{thm:g11:forces-and-motion:secondlaw}
\index{牛顿第二定律}
合力 $\vect F$ 作用于质量为 $m$ 的物体，作用很短的时间，
$\Delta t$ 内，其质心速度改变一个矢量
$\Delta \vect v$，其\emph{方向}与 $\vect F$ 相同，且其大小
与 $F$、$\Delta t$ 成正比，
与质量成反比：
\[
\Delta v \propto \frac{F \, \Delta t}{m} .
\]
\end{theorem}
\admitted

\begin{remark}
力或作用时间加倍，速度改变量加倍；质量加倍，
速度改变量减半。下一学年将把它精确化为方程，而本科一年级分册
会探讨它在哪些参考系中成立；今年用比例关系就已足够——
它们已经用于安全气囊、刹车和轨道之中。
\end{remark}

\begin{proposition}[三种典型情形]\label{prop:g11:forces-and-motion:cases}
设物体以速度 $\vect v$ 运动，受合力 $\vect F$ 作用。
\begin{enumerate}
  \item $\vect F$ 与 $\vect v$ 同向：物体沿直线\emph{加速}；
  \item $\vect F$ 与 $\vect v$ 反向：物体沿直线\emph{减速}；
  \item $\vect F$ 与 $\vect v$ 垂直：物体\emph{转弯}，
        速率不变。
\end{enumerate}
\end{proposition}

\begin{proof}
将 $\Delta\vect v$（方向沿 $\vect F$，
见\cref{thm:g11:forces-and-motion:secondlaw}），首尾相接地加到
$\vect v$ 上：沿 $\vect v$ 方向时箭头依次相接而变长；反向时，它们
部分抵消而变短；垂直时，每个小的 $\Delta \vect v$
只使箭头倾斜而不改变长度——只有方向在转动。
\end{proof}

\begin{omfigure}
\begin{tikzpicture}[scale=0.9, font=\small]
  \draw[-Stealth, omDef, thick] (0,1.5) -- (1.4,1.5) node[above] {$\vect v$};
  \draw[-Stealth, omThm, very thick] (0,0.75) -- (0.8,0.75) node[above] {$\vect F$};
  \draw[-Stealth, omDef, thick] (0,0) -- (2.2,0) node[above] {$\vect v'$};
  \node at (1.1,-0.7) {加速};
\end{tikzpicture}
\qquad
\begin{tikzpicture}[scale=0.9, font=\small]
  \draw[-Stealth, omDef, thick] (0,1.5) -- (2.2,1.5) node[above] {$\vect v$};
  \draw[-Stealth, omThm, very thick] (0.8,0.75) -- (0,0.75) node[above] {$\vect F$};
  \draw[-Stealth, omDef, thick] (0,0) -- (1.4,0) node[above] {$\vect v'$};
  \node at (1.1,-0.7) {减速};
\end{tikzpicture}
\qquad
\begin{tikzpicture}[scale=0.9, font=\small]
  \draw[-Stealth, omDef, thick] (0,1.5) -- (2.0,1.5) node[above] {$\vect v$};
  \draw[-Stealth, omThm, very thick] (1.0,0.4) -- (1.0,1.2) node[right] {$\vect F$};
  \draw[-Stealth, omDef, thick] (0,-0.5) -- (30:2.0)
    node[above] {$\vect v'$};
  \node at (1.1,-1.1) {转弯};
\end{tikzpicture}

{\small 三种典型情形：合力拖动速度矢量的箭头
使其各自变长、变短，或转弯。}
\end{omfigure}

\begin{example}[侧向一击]\label{ex:g11:forces-and-motion:kick}
一个冰球以 \qty{6.0}{m/s} 向东滑行；一次短暂的侧向撞击附加
$\Delta\vect v = \qty{8.0}{m/s}$，方向正北。新速度为
原速度与之的矢量和：大小 $\sqrt{6.0^2 + 8.0^2} = \qty{10.0}{m/s}$，
方向 $\tan\theta = 8.0/6.0$，即北偏东 $\theta \approx 53^\circ$——
原速度并未抹去，变化量只是叠加其上。
\end{example}

\section{比较变化：力、时间、质量}

\begin{method}[比值推理]\label{met:g11:forces-and-motion:ratio}
比较两种情形时，不必解方程，而应列出比值：
速度变化量之比等于 $F$ 之比，乘以 $\Delta t$ 之比，
再除以 $m$ 之比。然后读出哪一个因素发生了变化，变化了多少。
\end{method}

\begin{example}[装载的购物车]\label{ex:g11:forces-and-motion:cart}
同样的推力作用同样的一秒，使空的
\qty{20}{kg} 购物车以 \qty{1.2}{m/s} 的速度驶出。装载至
\qty{60}{kg}——质量变为三倍，$F\Delta t$ 不变——它以
原速度的三分之一，即 \qty{0.40}{m/s}。质量即物体的惯性。
\end{example}

\begin{remark}[碰撞、安全气囊、屈膝]\label{rem:g11:forces-and-motion:airbag}
碰撞中，$\Delta v$ 与 $m$ 不可更改——乘员的
速率必须降为零——故 $F\,\Delta t$ 为定值。唯一可调的
余地是\emph{时间}：延长 $\Delta t$，力便按同一比例减小。
头部若在安全气囊上停下，用时 \qty{0.08}{s}，而非在
方向盘上以 \qty{0.01}{s} 停下，所受平均力仅为八分之一；
吸能车身、头盔泡沫和屈膝也起同样的作用。
\end{remark}

\section{摩擦力：阻碍滑动的力}

\begin{definition}[静摩擦力与动摩擦力]\label{def:g11:forces-and-motion:friction}
相互接触的表面间的摩擦力，阻碍它们之间的\emph{相对
滑动}。\emph{静摩擦力}\index{摩擦力!静摩擦力}在二者不滑动时起作用，
大小自动调节至最大值，以抵消任何企图开始滑动的力；
\emph{动摩擦力}\index{摩擦力!动摩擦力}，
一旦滑动开始，其大小近似恒定，方向与滑动方向相反，
通常比静摩擦力的最大值\emph{更小}。
\end{definition}

\begin{example}[推衣柜]\label{ex:g11:forces-and-motion:wardrobe}
轻推一个衣柜：它不动——静摩擦力抵消了
你施加的 \qty{150}{N} 的力。再用力推：用力达到 \qty{400}{N} 时，它猛地开始
运动，此后用较小的力即可使其滑动——较弱的动摩擦力
便起了主导作用。
\end{example}

\begin{remark}[ABS 制动为何有效]\label{rem:g11:forces-and-motion:abs}
滚动车轮与路面的接触点不发生滑动：轮胎
依靠\emph{静摩擦力}抓地。抱死的车轮会打滑，只能获得
较弱的\emph{动摩擦力}，其方向始终与打滑方向相反，
无论车轮如何运动——转向失效。ABS 会在车轮抱死的瞬间
松开刹车：抓地力更强，制动距离更短，驾驶员仍能转向。
\end{remark}

\begin{omfigure}
\begin{tikzpicture}
\begin{axis}[omaxis, width=8.6cm, height=5.2cm, ybar, bar width=16pt,
  symbolic x coords={dry road, wet road, packed snow}, xtick=data, xticklabels={干燥路面，湿滑路面，压实积雪},
  xlabel={}, ylabel={制动距离（\unit{m}）}, ymin=0, ymax=185,
  nodes near coords,
  every node near coord/.append style={font=\small},
  enlarge x limits=0.25]
\addplot[draw=omDef, fill=omDef!25] coordinates
  {(dry road,40) (wet road,70) (packed snow,160)};
\end{axis}
\end{tikzpicture}

{\small 从 \qty{90}{km/h} 出发的典型制动距离：轮胎与路面之间的
摩擦力是使汽车减速的唯一力，积水或积雪会大幅削弱它。}
\end{omfigure}

\section{转弯：圆周运动需要指向圆心的力}

\begin{definition}[匀速圆周运动]\label{def:g11:forces-and-motion:ucm}
物体处于\emph{匀速圆周运动}\index{匀速圆周运动}，
即其轨迹是以恒定速率绕行的圆时。速度沿圆的切线方向，
每一时刻方向都在改变——
因此速度从不恒定。
\end{definition}

\begin{proposition}[指向圆心的力]\label{prop:g11:forces-and-motion:inward}
做匀速圆周运动的物体所受合力
始终不为零：每一时刻都由物体指向
圆心。
\end{proposition}

\begin{proof}
速率恒定，故合力沿 $\vect v$ 方向无分量
（情形 1 与 2，见\cref{prop:g11:forces-and-motion:cases}）：合力
垂直于切线，因而沿半径方向——且指向圆心，
因为每一个 $\Delta \vect v$ 都使路径弯向圆心。
\end{proof}

\begin{omfigure}
\begin{tikzpicture}[scale=1.0, font=\small]
  \draw[gray!50, dashed] (0,0) circle (1.7);
  \fill (0,0) circle (1.5pt) node[below] {手};
  \draw[gray!70] (0,0) -- (90:1.7) (0,0) -- (210:1.7) (0,0) -- (330:1.7);
  \fill[omExo] (90:1.7) circle (2.5pt);
  \fill[omExo] (210:1.7) circle (2.5pt);
  \fill[omExo] (330:1.7) circle (2.5pt);
  \draw[-Stealth, omThm, very thick] (90:1.7) -- (90:0.75) node[pos=0.35, right] {$\vect F$};
  \draw[-Stealth, omThm, very thick] (210:1.7) -- (210:0.75);
  \draw[-Stealth, omThm, very thick] (330:1.7) -- (330:0.75);
  \draw[-Stealth, omDef, thick] (90:1.7) -- ++(180:1.2) node[left] {$\vect v$};
  \draw[-Stealth, omDef, thick] (210:1.7) -- ++(300:1.2);
  \draw[-Stealth, omDef, thick] (330:1.7) -- ++(60:1.2);
\end{tikzpicture}

{\small 用绳子甩动的小球：速度沿切线方向，拉力
指向圆心。剪断绳子：小球沿切线方向飞出。}
\end{omfigure}

\begin{example}[绳子、弯道与月球]\label{ex:g11:forces-and-motion:moon}
每一次匀速转动背后都隐藏着指向内侧的力。甩动的小球：绳子的
张力。汽车过弯道：路面作用于轮胎的侧向静摩擦力——冰面上若没有这一摩擦力，
汽车便沿直线冲出。月球：地球的
万有引力（\cref{ch:g10:universal-gravitation}），使它
在 \num{27.3} 天内绕地球转一周——它不是被\emph{托}住，而是被侧向拉入
永不停息的转动之中，绕着地球不断“下落”；
下学年将它发展为卫星力学。
\end{example}

\section{练习}

\begin{exercise}[$\star$]\label{exo:g11:forces-and-motion:1}
判断下列情形中合力是否为零，并由速度说明理由：（a）汽车以
\qty{90}{km/h} 的速率匀速直行；（b）汽车沿直线刹车；
（c）汽车以 \qty{50}{km/h} 的速率匀速过弯道；（d）跳伞者
以终端速度下落。
\end{exercise}

\begin{exercise}[$\star$]\label{exo:g11:forces-and-motion:2}
质量分别为 \qty{3.0}{kg} 和 \qty{1.0}{kg} 的两个小球，位于
长 \qty{0.80}{m} 的轻杆两端。求质心位置；若两球质量相等，质心在何处？
\end{exercise}

\begin{exercise}[$\star$]\label{exo:g11:forces-and-motion:3}
同样的短暂推力分别作用于空的 \qty{15}{kg} 行李车和
装载至 \qty{45}{kg} 的同一辆行李车上。比较两者的速度
变化；要使满载行李车获得与空车相同的 $\Delta v$，需要多大的推力？
\end{exercise}

\begin{exercise}[$\star$]\label{exo:g11:forces-and-motion:4}
一辆有轨电车以 \qty{10}{m/s} 的速率刹车；一秒后
速率变为 \qty{8}{m/s}，方向不变。求 $\Delta \vect v$ 及合力的方向。
这属于\cref{prop:g11:forces-and-motion:cases}中的哪一种情形？
\end{exercise}

\begin{exercise}[$\star$]\label{exo:g11:forces-and-motion:5}
判断是静摩擦力还是动摩擦力——并说明方向：（a）箱子抵抗你（过轻的）
推力；（b）同一箱子开始滑动后；
（c）书放在略微倾斜的托盘上，未滑动。
\end{exercise}

\begin{exercise}[$\star\star$]\label{exo:g11:forces-and-motion:6}
合外力 $F$ 作用时间 $\Delta t$ 于质量为 $m$ 的物体，
产生 $\Delta v = \qty{4.0}{m/s}$。预测下列情形的 $\Delta v$：（a）$2F, \Delta t, m$；
（b）$F, \Delta t/2, m$；（c）$F, \Delta t, 2m$；（d）$3F, 2\Delta t, 6m$。
\end{exercise}

\begin{exercise}[$\star\star$]\label{exo:g11:forces-and-motion:7}
冰球以 \qty{8.0}{m/s} 的速率向东滑动；一次击打使其
获得 $\Delta \vect v = \qty{6.0}{m/s}$、方向向北的速度变化，求新的速率和方向。
为什么速率不是 $8.0 + 6.0 = \qty{14}{m/s}$？
\end{exercise}

\begin{exercise}[$\star\star$]\label{exo:g11:forces-and-motion:8}
守门员以两种方式接住同一脚射门：手臂绷紧（球在
\qty{0.02}{s} 内停止），或随球“缓冲”（用时 \qty{0.08}{s}）。比较两种方式的平均作用力，
并说明她在选择之前已固定了哪些量。
\end{exercise}

\begin{exercise}[$\star\star$]\label{exo:g11:forces-and-motion:9}
起跑时，是哪一个外力推动短跑运动员向前？
为什么钉鞋有帮助——又为什么在
绝对光滑的冰面上，任何人都无法加速，甚至无法行走？
\end{exercise}

\begin{exercise}[$\star\star$]\label{exo:g11:forces-and-motion:10}
紧急刹车时，为什么装有 ABS（车轮保持滚动）的汽车
通常比车轮抱死、侧滑的汽车停得更近——而侧滑的汽车为什么
不能再被操纵？引用
\cref{def:g11:forces-and-motion:friction}。
\end{exercise}

\begin{exercise}[$\star\star$]\label{exo:g11:forces-and-motion:11}
链球运动员甩动连着铁丝的球。是什么力使球
转动？画出释放后球的轨迹。月球为什么不需要铁丝，
若没有引力，它的轨迹会变成什么？
\end{exercise}

\begin{exercise}[$\star\star\star$]\label{exo:g11:forces-and-motion:12}
扫帚可看作质量为 \qty{1.0}{kg} 的均匀细杆（其质心距顶端
\qty{0.60}{m}），顶端装有质量为 \qty{2.0}{kg} 的扫把头，其质心距顶端
\qty{1.20}{m}。将每一部分视为质点，
求扫帚的质心位置，并检验
\cref{def:g11:forces-and-motion:com}中的反比关系。
\end{exercise}

\begin{exercise}[$\star\star\star$]\label{exo:g11:forces-and-motion:13}
质量为 $m$ 的小车与质量为 $2m$ 的满载货车，以相同速率、
施加相同制动力刹车。
\begin{enumerate}
  \item 比较两者停下所需的时间。
  \item 两者速率均匀减小，因此以相同的
        \emph{平均}速率行驶：比较两者的停车距离。
  \item 小车现在以\emph{两倍}速率开始刹车，制动力相同：
        与第一次停车相比，比较其时间与距离。这对限速有何启示？
\end{enumerate}
\end{exercise}

\begin{exercise}[$\star\star\star$]\label{exo:g11:forces-and-motion:14}
月球绕地球运行，轨道半径 $R \approx \qty{3.84e8}{m}$，周期 $T = \num{27.3}$ 天。
\begin{enumerate}
  \item 求它的轨道速率（周长除以周期）。
  \item 它的速度变化方向指向何处？
        由哪个力提供？
  \item 利用\cref{prop:g11:forces-and-motion:inward}改进下述说法：“
        月球不会下落，是因为它运动得足够快”。
\end{enumerate}
\end{exercise}

\begin{exercise}[$\star\star\star$]\label{exo:g11:forces-and-motion:15}
一个质量为 \qty{50}{g} 的鸡蛋以 \qty{4.0}{m/s} 的速度撞击地面：若落在瓷砖上，
约经 \qty{1}{ms} 停下；若落在厚海绵上，约经 \qty{50}{ms} 停下。比较
两种情形下的平均作用力。用同样的道理解释，
跳跃落地时为什么要屈膝，以及体操运动员的缓冲垫由什么材料制成。
\end{exercise}

\section{问题：碰撞测试}

\begin{problem}[{周末问题——设计温和的碰撞：
既然碰撞决定了速度的变化量，汽车上的每一种安全装置
都是延长时间的机器}]
\label{pb:g11:forces-and-motion:1}
一个碰撞测试实验室让一辆汽车以 \qty{50}{km/h} 的速度撞向刚性墙，
车内载有一个 \qty{70}{kg} 的成年假人和一个 \qty{20}{kg} 的儿童假人；
摄像机记录每一次停止的时间。你的任务是：找出碰撞中的“暴力”藏在何处，
并通过设计将其消除。

\medskip
\noindent\textbf{第一部分——碰撞的度量。}
\begin{enumerate}
  \item 将 \qty{50}{km/h} 换算为 \unit{m/s}。
  \item 在任何正面停车中，乘员的速度变化量
        大小相同。这是哪一个量？为什么没有装置能改变它？
  \item 引用\cref{thm:g11:forces-and-motion:secondlaw}：当 $\Delta v$
        与 $m$ 固定时，$F\,\Delta t$ 也固定。设计者唯一可调的旋钮是什么？
  \item 墙壁使汽车在 \qty{0.050}{s} 内停下；吸能区
        将这一时间延长至 \qty{0.150}{s}。比较两者的平均作用力。
  \item 在停车过程中，$\Delta \vect v$ 与假人所受合力的方向
        分别指向哪个方向？
\end{enumerate}

\medskip
\noindent\textbf{第二部分——吸能区。}
\begin{enumerate}[resume]
  \item 用一句话回答：工程师为何要把汽车前部
        设计成在碰撞中被毁坏？
  \item 在一辆老式刚性汽车中（在 \qty{0.050}{s} 内停下），假人被
        刚性地绑在座椅上。比较其所受作用力与现代汽车 \qty{0.150}{s} 停车时的作用力。
  \item 停车过程中，汽车的速率均匀下降（平均值为初速率的一半）。
        它在 \qty{0.150}{s} 内行驶了多远？与真实吸能区相比如何？
  \item 为何不把吸能区做到 \qty{10}{m}？车上哪一部分
        必须\emph{不}溃缩？
  \item 将第二部分总结为一句口号：“你无法选择 $\Delta v$，
        你必须选择\dots”。
\end{enumerate}

\medskip
\noindent\textbf{第三部分——安全带、安全气囊与儿童。}
\begin{enumerate}[resume]
  \item 未系安全带的假人保持其 \qty{14}{m/s} 的速度（见\cref{ch:g10:inertia}），
        直到挡风玻璃在 \qty{0.005}{s} 内将其挡停；系好安全带的假人
        则随汽车在 \qty{0.150}{s} 内停下。比较两者的平均作用力。
  \item 为什么略有伸长的安全带比
        不可伸长的钢制束带更好？
  \item 若头部撞在方向盘上，将在 \qty{0.010}{s} 内停下；安全气囊
        将这一时间延长至 \qty{0.080}{s}。比较两者的作用力。
  \item 成人与儿童经历相同的 $\Delta v$ 与相同的
        $\Delta t$。比较两条安全带必须提供的合力。
  \item 儿童坐在大人腿上：手臂必须使 \qty{10}{kg} 的婴儿
        在 \qty{0.150}{s} 内使速度改变 \qty{14}{m/s}，而重力
        使速度每秒改变 \qty{9.8}{m/s}。手臂需要快多少倍
        才行——托住多大质量的物体，才会感到同样的
        沉重？请作出结论。
  \item 儿童所受的力较小——那为何还需要儿童座椅？想一想
        安全带在\emph{何处}施力，以及 $\Delta t$。
\end{enumerate}

\medskip
\noindent\textbf{第四部分——安全笼与结论。}
\begin{enumerate}[resume]
  \item 乘员舱是刚性的，而两端会溃缩。为什么乘员舱
        必须不发生变形，尽管刚性意味着剧烈的停止？
  \item 两辆相同的汽车以各 \qty{50}{km/h} 的速度迎面相撞。
        它们各自在何处停下？将每次碰撞与墙壁测试相比较。
  \item 一辆小汽车与一辆满载的卡车迎面相撞。它们相互施加的力
        大小相等、方向相反，正如一切相互作用都是相互的
        （见\cref{ch:g11:fundamental-interactions}）。谁的速度
        变化更大？
  \item 列出本题中的各种安全装置，以及它们每一种
        所延长的同一个物理量。
  \item 总结：用一句话写出温和碰撞的设计规律，
        只许使用\emph{质量}、\emph{速度变化量}、
        \emph{力}与\emph{时间}。
\end{enumerate}
\end{problem}
